\chapter{Conclusion and Future Work}
\label{chap:conclusion}

% -----------------------------------------------------------------------------
\section{Conclusion}
\label{sec:conclusion}
% -----------------------------------------------------------------------------

This thesis asked whether a coordinated workflow of specialised LLM
agents could produce a clinically plausible differential diagnosis from
longitudinal patient evidence while exposing its reasoning to a human
reviewer. The Clinical Multi-Agent Decisioning System (CMADS) was
designed and evaluated as a direct response to that question. Its
architecture is described in Chapter~\ref{chap:methodology} and its
results are reported in Chapter~\ref{chap:results}.

The research question stated in Section~\ref{sec:aim} is answered in two
parts. On the first half --- whether a small team of specialised agents
can reach a clinically plausible diagnosis when they review, critique, and
remember each other's work --- the seven-agent pipeline produced the
headline rates reported in Section~\ref{sec:concl_findings}, and a
matched single-LLM baseline on the same 160-patient cohort
(Section~\ref{sec:results_single_llm}) attributes the bulk of those
rates to the orchestration rather than to the underlying model. On the
second half --- whether giving the team memory of past patients
actually helps --- the paired controlled comparison of 160 identical
patients run with and without the memory subsystem gives a clear
answer: multi-level memory wins by $+23.1$ percentage points on
DIRECT ($76.9\,\%$ vs.\ $53.8\,\%$, 123/160 vs.\ 86/160) and adds
$+1.9$ percentage points on Found rate
($94.4\,\%$ vs.\ $92.5\,\%$) on the same patients. The corresponding
hypothesis is therefore \emph{supported}: the non-destructive
Reviewer-Refiner pattern is the practical mechanism that lets the
system rewrite a flawed primary diagnosis while preserving the
earlier reasoning for audit, and multi-level memory delivers a large
gain on top-1 accuracy.

All nine objectives stated in Section~\ref{sec:objectives} were
implemented; they were evaluated where reported in
Chapter~\ref{chap:results}, within the bounds defined in
Section~\ref{sec:scope}. In Section~\ref{sec:objectives} order:
\textbf{(O1)} the clinical data pipeline (Bronze/Silver/Gold over
Synthea FHIR-R4);
\textbf{(O2)} the 160-patient evaluation cohort built by the
deterministic eligibility rules and the LLM-based detectability
verifier;
\textbf{(O3)} the seven-agent LangGraph multi-agent workflow;
\textbf{(O4)} the Reviewer-Refiner pattern that allows structured
clinical review while preserving the earlier diagnostic output for
audit;
\textbf{(O5)} the four-tier shared-memory subsystem covering working,
episodic, semantic, and case-based tiers;
\textbf{(O6)} the NICE-guideline-grounded Treatment Planning stage
gated on DIRECT matches;
\textbf{(O7)} the LLM-as-judge evaluation framework using
DIRECT~/~INDIRECT~/~MISS with Rank-1-within-found;
\textbf{(O8)} the paired 160-patient memory A/B and the single-LLM
controlled baseline against the multi-agent pipeline;
\textbf{(O9)} the React + FastAPI doctor console with its two journeys
(Doctor runtime and Researcher statistics) for clinician and reviewer
inspection.
A single-LLM controlled baseline on the same 160-patient cohort
(Section~\ref{sec:results_single_llm}) confirms that the multi-agent
orchestration buys $+33.8$ percentage points on DIRECT and
$+20.0$ percentage points on Found over a single-call baseline at
identical patient data and identical reasoning model, on a paired
comparison of the same 160 patients.
The literature gaps identified in Chapter~\ref{chap:background} --- end-to-end
coordination, non-destructive deliberation, persistent shared memory, and
leakage-aware evaluation --- are each addressed directly by the
corresponding methodology section, leaving adversarial robustness of
multi-agent topologies as the principal gap that this thesis did not
engage with empirically and that is returned to as future work.

% -----------------------------------------------------------------------------
\section{Contributions}
\label{sec:concl_contributions}
% -----------------------------------------------------------------------------

The substantive contributions of this thesis are:

\begin{enumerate}

  \item \textbf{A LangGraph clinical multi-agent pipeline.} A seven-agent
  workflow that compiles end-to-end coordination, parallel evidence
  extraction, adaptive diagnostic reasoning, and non-destructive
  Reviewer-Refiner deliberation into a single typed graph.

  \item \textbf{A four-tier shared-memory subsystem.} A clinical
  realisation of the CoALA memory taxonomy, with working, episodic,
  semantic, and case-based tiers exposed to agents through a single
  facade.

  \item \textbf{A reproducible Synthea-to-Gold data pipeline.} A medallion
  architecture (Bronze, Silver via OMOP CDM v5.4, Gold) gated by an
  LLM-based detectability verifier, producing per-patient case files with
  hidden ground-truth labels.

  \item \textbf{A NICE-guideline-grounded treatment layer.} Qdrant-backed
  retrieval that is invoked only on DIRECT-matched diagnoses, keeping
  guideline-aligned recommendations from being attached to uncertain
  primaries.

  \item \textbf{A transparent LLM-as-judge evaluation framework.}
  Three-label adjudication of free-text differentials, rank-1-when-found, and
  a paired controlled memory A/B on 160 identical patients, reported
  transparently as a controlled head-to-head rather than as an automatic
  accuracy gain.

  \item \textbf{A doctor-facing inspection dashboard.} A single React +
  FastAPI application with two journeys --- \emph{Doctor (runtime)} for a
  live patient run and \emph{Researcher (statistics)} for cohort-level
  audit of past runs --- exposing the agent chain, shared-memory
  operations, retrieved evidence, and the final ranked differential for
  clinical review.

\end{enumerate}

% -----------------------------------------------------------------------------
\section{System properties beyond accuracy}
\label{sec:concl_properties}
% -----------------------------------------------------------------------------

Match-type accuracy is the empirical headline of this thesis, but the
architecture was designed against four engineering properties that
matter for a clinical assistive system. Two of these
properties --- explainability and reproducibility --- are exhibited by
the running system and are reported as demonstrated. The remaining
two --- extendability and scalability --- are \emph{design
properties} that follow from the architectural choices made in
Chapter~\ref{chap:methodology}, but that this evaluation did not
empirically test at the scale that would settle them; they are
labelled accordingly below.

\paragraph{Explainability (demonstrated).}
Every agent's intermediate reasoning, the memory operations (which
prior case was retrieved, why, and how the diagnostic agent used it),
and every NICE-guideline passage that informed the treatment plan are
stored as typed JSON and exposed verbatim in a doctor-facing console.
A clinician can inspect why the system reached its conclusion at any
layer --- final differential, reviewer critique, retrieved prior,
guideline passage --- rather than trust an opaque box. The console is
the same artefact the researcher uses to debug a regression, which
keeps the explanation surface and the development surface coherent.

\paragraph{Extendability (design property; not empirically validated here).}
The architecture is designed to allow new evidence modalities to be
added without re-architecting the diagnostic chain. All heavy data
is moved into the Gold layer (one JSON per patient), agents are
treated as readers of that layer, and adding a new modality ---
X-ray, MRI, pathology slide, ECG strip, or free-text discharge
summary --- amounts in principle to writing one new converter agent
that lowers the modality into structured text fields appended to the
Gold-layer \texttt{ehr\_case.json}. This separation of concerns is
an architectural choice; no per-modality extension experiment was
conducted in this work, so the claim is restricted to design intent.
Demonstrating that the existing diagnostic, reviewer, refiner, and
treatment agents in fact consume an additional modality without
behavioural regression is left to future work.

\paragraph{Scalability (design property; not empirically validated here).}
The LLM adapter is provider-agnostic (Groq, OpenAI, Anthropic,
Gemini, Ollama); the runtime is designed to pick an engine from the
configured keys at request time, and the same pipeline is intended
to run against any of them without code changes. Local execution
via Ollama works on a single workstation, while cloud routing is
designed to scale to whatever the configured provider supports. The
doctor console's runtime picker is populated dynamically from the
live provider catalogue, so no specific model is hard-coded. No
per-modality or per-provider scaling experiment was conducted in
this work, however; the cross-provider robustness, latency, and
cost behaviour that would justify a scalability claim at the level
of additional providers and additional cohort sizes is left to future
work.

\paragraph{Reproducibility (demonstrated).}
Every pipeline run writes a typed execution trace
(\texttt{execution\_trace.json}) and a per-agent JSON output to disk.
Any number reported in this thesis can be re-derived by re-aggregating
those artefacts --- there are no hidden state stores or one-shot
computations. The same property makes the evaluation harness
deterministic: re-judging the cohort against a different ground-truth
ontology, or with a different judge model, requires only re-running
the evaluator stage against the saved differentials.

Together these four properties describe the design intent of CMADS:
explainability and reproducibility are exhibited by the running
system and supported by the on-disk traces, while extendability and
scalability are architectural choices whose empirical validation
across additional modalities, providers, and cohort scales is left
to future work.

% -----------------------------------------------------------------------------
\section{Key Empirical Findings}
\label{sec:concl_findings}
% -----------------------------------------------------------------------------

Three empirical findings of this thesis are worth highlighting at the
point of conclusion. First, within this synthetic-evaluation setting,
the end-to-end pipeline produced encouraging diagnostic results on
the 160-patient paired cohort with multi-level memory enabled:
a DIRECT match in $76.9\,\%$ of evaluated cases (123/160), a combined
DIRECT-or-INDIRECT found rate of $94.4\,\%$ (151/160), and the
ground-truth condition ranked first in $64.5\,\%$ of cases in which
it was found (Rank-1 within found), at a median per-patient latency
around two minutes and with no single patient run aborted by an
agent failure. The detailed per-disease-family breakdown and
per-agent reliability profile are reported in
Sections~\ref{sec:results_e2e} and \ref{sec:results_agents}.

Second, the paired controlled comparison of the principal
multi-level memory configuration against the single-level baseline
on the same 160 patients
(Section~\ref{sec:results_memory_ab}) gives a large and clean
effect. The DIRECT rate rises from \textbf{$53.8\,\%$} (86/160) for
the single-level baseline to \textbf{$76.9\,\%$} (123/160) for the
principal configuration, a $+23.1$ percentage-point advantage on the
same patients. The Found rate also rises, from $92.5\,\%$ to
$94.4\,\%$ ($+1.9$ pp), with a matching drop on MISS. The earlier
205-vs-100 unpaired aggregate that appeared in drafts of this work
is superseded: on the paired test that controls for cohort
composition, the full configuration's effect on DIRECT is large and
positive.

Third, the memory evaluation itself is a useful warning about how
memory-augmented clinical agents are commonly reported. The paired
controlled comparison in this thesis exposes a 205-vs-100 unpaired
result that earlier appeared to show a large DIRECT gap as a cohort
artifact. Reporting that pattern transparently --- paired comparison
as the primary result, unpaired numbers labelled as such --- is
important for any later work on memory-augmented clinical agents.

% -----------------------------------------------------------------------------
\section{Limitations}
\label{sec:concl_limitations}
% -----------------------------------------------------------------------------

Several limitations qualify these conclusions and are made explicit
here.

\paragraph{Synthetic data only.}
All patient records originate from the Synthea generator and have passed
the LLM-based detectability verifier of
Section~\ref{sec:llm_verifier}; the system's accuracy on real electronic
health records, which encode noisier evidence, more frequent missingness,
and the full distribution of human comorbidities, cannot be inferred
from the results reported in this thesis.

\paragraph{Small memory-comparison cohorts.}
The memory-comparison cohort is limited to $n = 160$ paired patients,
and the substantive finding --- that multi-level memory delivers a
$+23.1$ percentage-point DIRECT advantage and a complementary
$+1.9$ percentage-point Found gain over the single-level baseline on
identical patient records --- would benefit from a larger replication
across a more diverse disease distribution before the magnitude of the
effect can be treated as settled.

\paragraph{Cardiometabolic-only coverage and cohort imbalance.}
The eight target disease families covered in this thesis are dominated
by cardiometabolic conditions, with the ESRD~/~CKD family alone
accounting for almost half of the evaluated cohort. Acute presentations,
oncology, infectious disease, and paediatric cases are out of scope.
Synthea's disease state machines themselves are an upper bound on what
the evaluation can demonstrate; the system's behaviour on conditions
outside this distribution is not characterised by the results reported
here.

\paragraph{No clinician-in-the-loop validation.}
Diagnostic accuracy in this thesis is assessed by an LLM-as-judge
protocol against Synthea ground truth; no practising clinician
reviewed the outputs as part of the present study. Prospective
clinician review is identified as future work.

\paragraph{Single-model pipeline.}
The agent pipeline runs every stage on a single primary reasoning model
(GPT-OSS-120B), which is a simplification of what a production system
could do; the implications of routing different stages to different
models are returned to in Section~\ref{sec:concl_future}.

\paragraph{Single-judge LLM evaluator.}
All match-type adjudications are produced by a single LLM-as-judge
configuration (Qwen3-32B); replacing the judge with a different model or
ensembling multiple judges might shift the boundary between DIRECT and
INDIRECT classifications and is not investigated here.

The system remains a research artefact: it is positioned as an
assistive tool for a human reviewer, not an autonomous clinical
decision-maker.

% -----------------------------------------------------------------------------
\section{Future Work}
\label{sec:concl_future}
% -----------------------------------------------------------------------------

Five directions are most likely to extend the work reported in this
thesis.

\paragraph{Real-EHR validation.}
The most consequential next step is to repeat the evaluation on real
clinical data such as MIMIC-IV or eICU, and check whether the
architecture-over-scale and memory-trade-off findings of
Chapter~\ref{chap:results} replicate outside the synthetic setting.

\paragraph{Prospective clinician-in-the-loop study.}
A prospective multi-reader study, in which several clinicians review
CMADS outputs and matched single-LLM outputs on a shared cohort
blinded to the source, would isolate the contribution of the
multi-agent architecture to clinical usefulness as distinct from raw
match-type accuracy.

\paragraph{Multi-model routing across agents.}
The current pipeline uses a single primary reasoning model for all
seven agents. The agents have different cognitive requirements
nonetheless --- structured extraction, open-ended reasoning,
adversarial critique, and retrieval-grounded application --- and
assigning each role to a capability- and cost-matched model, with
dynamic routing by case complexity in the spirit of
MDAgents~\cite{Kim2024MDAgents} and TAO~\cite{Kim2025TAO}, would
concentrate compute on the steps that benefit most. The
provider-agnostic adapter already supports this as a configuration
change.

\paragraph{Guideline coverage beyond NICE.}
Extending the retrieval corpus beyond NICE --- to NHS, WHO, ESC, and
specialty-society guidelines --- would let the agent reason under more
than one standard of care and flag inter-guideline disagreement as an
additional signal to the reviewing clinician.

\paragraph{Memory-architecture ablations.}
The memory finding of Chapter~\ref{chap:results} deserves a dedicated
follow-up. A tier-by-tier ablation should run each memory configuration on the
same patient cohort and keep the case store population separate from the test
cohort. That would identify which tier drives the differential-broadening
behaviour and whether the precision trade-off can be recovered by selective
gating. The system records every memory operation in a typed trace, so the
ablation is straightforward to instrument: one-tier-out runs on the same
160-patient cohort would isolate the contribution of case-based retrieval,
episodic summaries, and the combination.

% -----------------------------------------------------------------------------
\section{Closing Remarks}
\label{sec:concl_remarks}
% -----------------------------------------------------------------------------

CMADS is an assistive system, not a replacement for a clinician. The
seven-agent pipeline, the four-tier shared memory, and the doctor
console are designed so that the human reviewer remains the decision
authority and can inspect every step the agents took to reach a
proposed diagnosis. The takeaway of the thesis is more modest than a
headline accuracy number: a careful coordination architecture and a
leakage-aware evaluation protocol matter as much as the model itself
when building clinical multi-agent systems that are meant to be
reviewed rather than trusted blindly.
